\paragraph{Class LoRA.} The adapter was trained on 16 training patches
for one epoch with the classification target and the full Condition~C
card (final training loss 1.84), and tested on the full test set. This
is a small-scale configuration, not a complete fine-tune of the 256
training patches. Exploratory runs on all 256 training patches for two
to three epochs reached 100\% development and test accuracy, because
the model learned to copy \texttt{candidate\_class} from the summary
line. These runs are excluded from the results.

\paragraph{Explanation LoRA.} Qwen2-VL-2B itself acted as teacher: it
generated a justification for each training patch, and a generation was
kept only if its class matched the label, it contained at least two
checkable claims, and all of those claims agreed with the measured
changes. Of 256 training patches, 28 passed. A larger 7B teacher could
not be used because its download exceeded the cluster's disk quota. The
training run was capped at the first 128 training patches, so 16 of the
28 filtered justifications were used. Both training and evaluation used
the card without the label, the checkpoint was selected on the
development set, and the test set was evaluated in explanation mode.

\subsubsection{Evaluation Metrics}

Classification is scored by accuracy with Wilson 95\% intervals, macro-F1, and comparison with the 34.5\%
majority-class baseline; McNemar's test  compares
paired predictions.

Explanations are scored automatically. The scorer searches the text for
an indicator (NDVI, EVI, NBR, NDMI, rainfall, temperature, soil
moisture, VV, VH) and pairs it with the nearest direction word
(increase-type, decrease-type or stable-type). Each measured change is
converted to a direction using an indicator-specific dead-band, below
which a change counts as stable. A claim is \emph{correct} when its
direction matches the measured direction. We report the mean number of
checkable claims per explanation, the share of correct claims (pooled
over patches), the share of explanations with at least one incorrect
claim, and self-consistency (whether the dominant claimed direction
agrees with the model's own class). ROUGE-L  is
computed but not used as a measure of faithfulness, because overlap with
the supplied text rewards copying rather than reasoning. The scorer measures consistency with the evidence,
not faithfulness to the model's internal reasoning.

\subsection{Experimental Workflow}

Figure~\ref{fig:workflow} summarises the pipeline.

\begin{figure}[t]
\centering
\includegraphics[width=\columnwidth]{fig_workflow.png}
\caption{EcoExplain-VLM workflow. After the dataset is frozen, three
experiments are run on the same 55 test patches and are scored
separately.}
\label{fig:workflow}
\end{figure}

The work proceeded in six stages.

\begin{enumerate}
\item \textbf{Study design.} The study area is Kanha Tiger Reserve.
Matched dry-season windows are T0 (January--February 2022) and T1
(January--February 2023).

\item \textbf{Collect satellite and environmental data.} Sentinel-2
(optical), Sentinel-1 (radar), CHIRPS (rainfall), ERA5-Land (climate)
and NASADEM (terrain) were selected as the five sources.

\item \textbf{Google Earth Engine.} Composites were built for each
window and 15 GeoTIFF files were exported for every complete patch
(366 patches; 5,490 files).

\item \textbf{Process data and create labels.} DIP produced the RGB
chips and tensors. Patch-mean NDVI change determined the automatic
label (\texttt{candidate\_class}). Spectral and environmental values
were written into a structured evidence card. Each patch became one
JSONL record: images + evidence text + \texttt{candidate\_class}.

\item \textbf{Freeze the dataset split.} Records were split
256/55/55 (train/development/test) by stratified sampling on the
label (seed 42; 70/15/15).

\item \textbf{Run three separate experiments on the same 55 test
patches.}
\begin{itemize}
\item \textbf{A. Classification.} Four VLMs received the images and
the full evidence card (Condition~C). The run was repeated after
removing the summary line that contains the label
(\texttt{evidence\_only}). The measure is class accuracy.
Supporting analyses on this track were the closed-set probe and the
$\Delta$NDVI baselines.
\item \textbf{B. Class LoRA.} Qwen2-VL-2B was fine-tuned on 16
training patches for one epoch with the classification target, then
tested on the same 55 patches. The measure is class accuracy.
\item \textbf{C. Explanations.} The label was removed from the card
and models were asked to justify a prediction. Free-form prompts,
structured cite-only prompts and the explanation LoRA were compared.
The measures are claim correctness and class accuracy.
\end{itemize}
\end{enumerate}

The development set was used only for prompt wording, the closed-set
bias and the explanation LoRA checkpoint. No headline result was
chosen by looking at the test set.
% ============================================================
% 3. RESULTS AND DISCUSSION
% ============================================================

\section{Results and Discussion}

\subsection{Dataset Characteristics}

Table~\ref{tab:data} summarises the benchmark. The optical index changes
move together: $\Delta$NDVI correlates with $\Delta$NBR, $\Delta$NDMI
and $\Delta$EVI at 0.80, 0.77 and 0.78. Mean window rainfall fell from
69.1\,mm in 2022 to 25.1\,mm in 2023, the temperature anomaly rose from
$-1.52^\circ$C to $+0.88^\circ$C, and the soil-moisture anomaly fell
from $+0.024$ to $-0.043$. Mean VV and VH backscatter changed by $-0.47$
and $-0.48$\,dB. On average, therefore, the second dry season was drier
and slightly less green, which is why ``increase'' is a relative rank
and why many increase patches have $\Delta$ close to zero.

\begin{table}[t]
\centering
\small
\begin{tabular}{lr}
\toprule
Quantity & Value \\
\midrule
Study rectangle & 2,350.7\,km$^2$ \\
Candidate pool / final patches & 1,200 / 366 \\
Decline / stable / increase & 120 / 124 / 122 \\
Train / development / test & 256 / 55 / 55 \\
Test decline / stable / increase & 18 / 18 / 19 \\
Majority baseline (test) & 34.5\% \\
Human agreement & 8/19 (42\%) \\
Mean $\Delta$NDVI (s.d.) & $-0.016$ (0.027) \\
Patches with $|\Delta\mathrm{NDVI}| \ge 0.05$ & 34 \\
Card trend words disagree with label & 108 \\
\bottomrule
\end{tabular}
\caption{EcoExplain-VLM dataset summary. Labels are NDVI-change ranks
within this region and pair of years.}
\label{tab:data}
\end{table}

\subsection{Zero-Shot Model Performance}

Table~\ref{tab:zeroshot} reports zero-shot accuracy on the test set.
Qwen2-VL-2B is the only model above the majority baseline, at 56.4\%
(31/55; 95\% interval 43.3--68.6\%). Even this margin is not
statistically secure at $n = 55$: McNemar's test against always
predicting ``increase'' gives $p = 0.12$. The other three models do not
separate the classes. LLaVA-Interleave-0.5B answered ``decrease'' for
all 55 patches, PaliGemma answered ``decline'' for 42 of 55, and BLIP-2
repeated its own prompt, from which the parser read the first class word
(``decline''). BLIP-2's 32.7\% is therefore a parsing artefact rather
than a prediction.

\begin{table*}[t]
\centering
\small
\begin{tabular}{llrcr}
\toprule
Model or baseline & Input & Acc. & 95\% CI & Macro-F1 \\
\midrule
Qwen2-VL-2B & images + full card & \textbf{56.4\%} & 43.3--68.6 & 0.466 \\
PaliGemma-3B & images + NDVI line & 36.4\% & 24.9--49.6 & 0.271 \\
LLaVA-Interleave-0.5B & images + full card & 32.7\% & 21.8--45.9 & 0.164 \\
BLIP-2 OPT-2.7B & stitched image only & 32.7\% & 21.8--45.9 & 0.164 \\
\midrule
LLaVA-Interleave-0.5B, \texttt{evidence\_only} & images + card without label line & 36.4\% & 24.9--49.6 & 0.251 \\
LLaVA-Interleave-0.5B, closed-set probe & class-word likelihood, dev-fitted bias & 60.0\% & 46.8--71.9 & 0.479 \\
\midrule
Majority class (always ``increase'') & --- & 34.5\% & 23.4--47.7 & 0.171 \\
$\Delta$NDVI rule, $\pm 0.02$ cut-offs & card $\Delta$NDVI & 63.6\% & 50.4--75.1 & 0.606 \\
$\Delta$NDVI rule, cut-offs fitted on train & card $\Delta$NDVI & 94.5\% & 85.1--98.1 & 0.945 \\
\bottomrule
\end{tabular}
\caption{Classification on the test set ($n = 55$). The top block is
generative zero-shot Condition~C. The middle block contains supporting
runs that are not zero-shot generative results. The bottom block uses no
model.}
\label{tab:zeroshot}
\end{table*}
